(1) 令 \(\alpha=\sqrt{\hbar/(2m\omega)}\)，位置算符是 \(x=\alpha(a+a^\dagger)\)，满足 \([a,a^\dagger]=1\)、\(a|n\rangle=\sqrt n|n-1\rangle\)。无扰动哈密顿量 \(H_0=\hbar\omega(a^\dagger a+1/2)\)，能级 \(E_n^{(0)}=\hbar\omega(n+1/2)\)。

(2) 由对易关系可得
\[
(a+a^\dagger)^2=a^2+a^{\dagger2}+2N+1.
\]
在 n 态计算四次期望值，可写成 \(\langle n|[(a+a^\dagger)^2]^2|n\rangle\)。因 \(a^2\) 和 \(a^{\dagger2}\) 改变量子数 \(\pm 2\)，只有总改变量零的项留下：常数对角贡献 \((2n+1)^2\)，以及 \(a^2a^{\dagger2}\) 与 \(a^{\dagger2}a^2\) 的期望分别为 \((n+1)(n+2)\) 与 \(n(n-1)\)。相加得
\[
\langle n|x^4|n\rangle
=\alpha^4[6n^2+6n+3]
=3\alpha^4(2n^2+2n+1).
\]

(3) 非简并微扰论一阶能级修正 \(\Delta E_n^{(1)}=\lambda\langle n|x^4|n\rangle\)，因此
\[
E_n\simeq\hbar\omega(n+1/2)
+3\lambda\left(\frac{\hbar}{2m\omega}\right)^2(2n^2+2n+1).
\]
基态修正为 \(3\lambda\alpha^4\)，第一激发态为 \(15\lambda\alpha^4\)，所以第一能隙
\[
E_1-E_0\simeq\hbar\omega+12\lambda\alpha^4.
\]
正四次势使高激发态增加的能量更多，能隙相应上升。

(4) 因 \(x^{4}\) 是偶宇称算符，它只能连接同宇称谐振子态；将 \((a+a^\dagger)^4\) 展开并正规序化，非对角矩阵元只可能有 \(\Delta n=\pm2,\pm4\)，还包含 \(\Delta n=0\) 的对角元，而奇数变化严格为零。一阶能级修正全部为正，因为 \(\lambda >0\) 且 \(x^4\ge0\) 为正算符；扰动理论的有效性要求修正远小于相邻无扰动能级间隔。
